\documentclass[10pt,a4paper]{amsart}
\usepackage{amsmath,amssymb,amsthm,graphicx}
% For commutative diagrams we use Xy-pic
\usepackage[all]{xy}

% We use 2cell for 2-commutative diagrams.
\xyoption{2cell}
% Theorem environments.
%
\theoremstyle{plain}
\newtheorem{theorem}[subsection]{Theorem}
\newtheorem{proposition}[subsection]{Proposition}
\newtheorem{lemma}[subsection]{Lemma}

\theoremstyle{definition}
\newtheorem{definition}[subsection]{Definition}
\newtheorem{example}[subsection]{Example}
\newtheorem{exercise}[subsection]{Exercise}
\newtheorem{situation}[subsection]{Situation}

\theoremstyle{remark}
\newtheorem{remark}[subsection]{Remark}
\newtheorem{remarks}[subsection]{Remarks}

\numberwithin{equation}{subsection}

% Macros
%
\def\lim{\mathop{\mathrm{lim}}\nolimits}
\def\colim{\mathop{\mathrm{colim}}\nolimits}
\def\Spec{\mathop{\mathrm{Spec}}}
\def\Hom{\mathop{\mathrm{Hom}}\nolimits}
\def\Ext{\mathop{\mathrm{Ext}}\nolimits}
\def\SheafHom{\mathop{\mathcal{H}\!\mathit{om}}\nolimits}
\def\SheafExt{\mathop{\mathcal{E}\!\mathit{xt}}\nolimits}
\def\Sch{\mathit{Sch}}
\def\Mor{\mathop{\mathrm{Mor}}\nolimits}
\def\Ob{\mathop{\mathrm{Ob}}\nolimits}
\def\Sh{\mathop{\mathit{Sh}}\nolimits}
\def\NL{\mathop{N\!L}\nolimits}
\def\CH{\mathop{\mathrm{CH}}\nolimits}
\def\proetale{{pro\text{-}\acute{e}tale}}
\def\etale{{\acute{e}tale}}
\def\QCoh{\mathit{QCoh}}
\def\Ker{\mathop{\mathrm{Ker}}}
\def\Im{\mathop{\mathrm{Im}}}
\def\Coker{\mathop{\mathrm{Coker}}}
\def\Coim{\mathop{\mathrm{Coim}}}

% Boxtimes
%
\DeclareMathSymbol{\boxtimes}{\mathbin}{AMSa}{"02}

%
% Macros for moduli stacks/spaces
%
\def\QCohstack{\mathcal{QC}\!\mathit{oh}}
\def\Cohstack{\mathcal{C}\!\mathit{oh}}
\def\Spacesstack{\mathcal{S}\!\mathit{paces}}
\def\Quotfunctor{\mathrm{Quot}}
\def\Hilbfunctor{\mathrm{Hilb}}
\def\Curvesstack{\mathcal{C}\!\mathit{urves}}
\def\Polarizedstack{\mathcal{P}\!\mathit{olarized}}
\def\Complexesstack{\mathcal{C}\!\mathit{omplexes}}
% \Pic is the operator that assigns to X its picard group, usage \Pic(X)
% \Picardstack_{X/B} denotes the Picard stack of X over B
% \Picardfunctor_{X/B} denotes the Picard functor of X over B
\def\Pic{\mathop{\mathrm{Pic}}\nolimits}
\def\Picardstack{\mathcal{P}\!\mathit{ic}}
\def\Picardfunctor{\mathrm{Pic}}
\def\Deformationcategory{\mathcal{D}\!\mathit{ef}}
\title{Existence of the second Adams operation on K\_0 of the vector bundles of a scheme, with additivity, pullback compatibility and multiplicativity}
\author{The Stacks Project Authors}
\date{stacks-project:chow.tex:local-snapshot-20261012; GNU FDL 1.2 or later}
\begin{document}
\maketitle
\noindent\textit{Editorial note.} All delivered mathematical text is the author's own, copied from the pinned Stacks Project chapter file chow.tex; no statement or proof step was rewritten, reordered or removed.\par\medskip
\noindent\textit{Editorial note.} Preamble provenance: the pinned chapter file opens with an input of the shared preamble and carries no preamble text of its own, so the project's single Stacks-wide preamble staged at sources/stacks-chow/preamble.tex is excerpted verbatim as two spans, lines 14-18 (Xy-pic setup) and lines 154-217 (the native theorem declarations and mathematical macros). Nothing was added, removed or redefined; the wrapper class and the amsmath/amssymb/amsthm/graphicx base packages are supplied by the assembler.\par\medskip
\noindent\textit{Editorial note.} The counted statement and proof contain no cross reference, no citation and no included file, so the delivered body is self contained apart from the notation of K\_0(Vect(X)) which the statement uses without recalling it.\par\medskip
\noindent\textit{Editorial note.} Canonical location of the counted proposition: Stacks tag 0FEJ (sources/stacks-chow/tags.txt line 19254).\par\medskip
\noindent\textit{Allowed standard prerequisites (editorial).}\begin{itemize}
\item The Grothendieck group K\_0 of finite locally free modules and its presentation by short exact sequences. Algebra, Section 'K-groups' (tag 00JC region) and the same chapter's earlier K\_0 material; the statement speaks of K\_0(Vect(X)) without recalling it. Supplies that K\_0(Vect(X)) is generated by classes of finite locally free modules subject to the relations coming from short exact sequences, which is the relation used at proof lines 12797-12800 to show that the formula psi\textasciicircum{}2 = [Sym\textasciicircum{}2] - [wedge\textasciicircum{}2] descends to a well defined additive map.
\end{itemize}\medskip
\begin{lemma}
\label{lemma-second-adams-operator}
Let $X$ be a scheme. There is a ring map
$$
\psi^2 :
K_0(\textit{Vect}(X))
\longrightarrow
K_0(\textit{Vect}(X))
$$
which sends $[\mathcal{L}]$ to $[\mathcal{L}^{\otimes 2}]$
when $\mathcal{L}$ is invertible and is compatible with pullbacks.
\end{lemma}

\begin{proof}
Let $X$ be a scheme.
Let $\mathcal{E}$ be a finite locally free $\mathcal{O}_X$-module.
We will consider the element
$$
\psi^2(\mathcal{E}) = [\text{Sym}^2(\mathcal{E})] - [\wedge^2(\mathcal{E})]
$$
of $K_0(\textit{Vect}(X))$.

\medskip\noindent
Let $X$ be a scheme and consider a short exact sequence
$$
0 \to \mathcal{E} \to \mathcal{F} \to \mathcal{G} \to 0
$$
of finite locally free $\mathcal{O}_X$-modules. Let us think of
this as a filtration on $\mathcal{F}$ with $2$ steps. The induced
filtration on $\text{Sym}^2(\mathcal{F})$ has $3$ steps with
graded pieces $\text{Sym}^2(\mathcal{E})$, $\mathcal{E} \otimes \mathcal{F}$,
and $\text{Sym}^2(\mathcal{G})$. Hence
$$
[\text{Sym}^2(\mathcal{F})] =
[\text{Sym}^2(\mathcal{E})] +
[\mathcal{E} \otimes \mathcal{F}] +
[\text{Sym}^2(\mathcal{G})]
$$
In exactly the same manner one shows that
$$
[\wedge^2(\mathcal{F})] =
[\wedge^2(\mathcal{E})] +
[\mathcal{E} \otimes \mathcal{F}] +
[\wedge^2(\mathcal{G})]
$$
Thus we see that
$\psi^2(\mathcal{F}) = \psi^2(\mathcal{E}) + \psi^2(\mathcal{G})$.
We conclude that we obtain a well defined additive map
$\psi^2 : K_0(\textit{Vect}(X)) \to K_0(\textit{Vect}(X))$.

\medskip\noindent
It is clear that this map commutes with pullbacks.

\medskip\noindent
We still have to show that $\psi^2$ is a ring map.
Let $X$ be a scheme and let $\mathcal{E}$ and $\mathcal{F}$
be finite locally free $\mathcal{O}_X$-modules.
Observe that there is a short exact sequence
$$
0 \to \wedge^2(\mathcal{E}) \otimes \wedge^2(\mathcal{F}) \to
\text{Sym}^2(\mathcal{E} \otimes \mathcal{F}) \to
\text{Sym}^2(\mathcal{E}) \otimes \text{Sym}^2(\mathcal{F}) \to 0
$$
where the first map sends $(e \wedge e') \otimes (f \wedge f')$ to
$(e \otimes f)(e' \otimes f') - (e' \otimes f)(e \otimes f')$ and
the second map sends $(e \otimes f) (e' \otimes f')$ to $ee' \otimes ff'$.
Similarly, there is a short exact sequence
$$
0 \to \text{Sym}^2(\mathcal{E}) \otimes \wedge^2(\mathcal{F}) \to
\wedge^2(\mathcal{E} \otimes \mathcal{F}) \to
\wedge^2(\mathcal{E}) \otimes \text{Sym}^2(\mathcal{F}) \to 0
$$
where the first map sends $e e' \otimes f \wedge f'$ to
$(e \otimes f) \wedge (e' \otimes f') + (e' \otimes f) \wedge (e \otimes f')$
and the second map sends
$(e \otimes f) \wedge (e' \otimes f')$ to
$(e \wedge e') \otimes (f f')$.
As above this proves the map $\psi^2$ is multiplicative.
Since it is clear that $\psi^2(1) = 1$ this concludes the proof.
\end{proof}

\end{document}
